% Summary of AI use for the developer documentation (Dean's Instruction No. I/2026), shown when \aideclarationtrue.
\section{Use of AI tools in this work}
\label{sec:ai-use}

As the faculty's AI policy requires, this section summarises how AI tools were used across the whole work; the signed
Declaration on the use of AI tools lists them in detail.

The program was developed with an AI coding assistant, Claude by Anthropic, used through Claude Code. I wrote the
requirements (Chapter~\ref{ch:spec}), chose the architecture and the technologies, and decided the important design
questions: that the language model may only propose a discount and deterministic code decides
(Section~\ref{sec:pricing}), that stock is locked in a fixed order and recorded in a ledger
(Section~\ref{sec:orders}), and that customers may reserve and pay at the stall. For each change I described what was
needed; the assistant wrote the code, the tests and the migrations; I reviewed the change, ran the test suites, and
tried it in the browser and on a phone before accepting it. Continuous integration (Section~\ref{sec:running}) ran every
suite on every change, and the evaluation (Section~\ref{sec:evaluation}) and the browser journey checked the system as a
whole. Several defects that the assistant introduced or missed were found this way and corrected, for example tests that
passed by chance and an option shown as sold out while the same colour was in stock.

The thesis text was drafted and edited with the same assistant from my outline, my notes and the code; I checked every
statement and number against the program and the test results and decided the structure and the content. The Russian
and Kyrgyz texts of the interface were drafted with it and corrected by me. The language model inside the product
(Section~\ref{sec:assistant}) is a component of the system, not a tool used to write it.

